\documentclass[11pt]{article}

\usepackage[margin=1in]{geometry}
\usepackage{amsmath,amssymb,mathtools,bm}
\usepackage{booktabs}
\usepackage{algorithm}
\usepackage{algpseudocode}
\usepackage{xcolor}
\usepackage{microtype}
\usepackage[hidelinks]{hyperref}

\newcommand{\Df}{\mathcal{D}_{\mathrm{f}}}
\newcommand{\Dr}{\mathcal{D}_{\mathrm{r}}}
\newcommand{\V}{\mathcal{V}}
\newcommand{\A}{\mathcal{A}}
\newcommand{\KL}{D_{\mathrm{KL}}}
\newcommand{\Ber}{\operatorname{Bern}}
\newcommand{\relu}[1]{\left[#1\right]_{+}}
\newcommand{\stopgrad}{\operatorname{sg}}
\newcommand{\FCP}{\textsc{FCP-MDLM}}

\title{Training Algorithms for Fact-wise Counterfactual Projection in Masked Diffusion Language Models}
\author{}
\date{}

\begin{document}
\maketitle

\section{Problem Setup}

Let $\theta_0$ denote the original masked diffusion language model (MDLM) trained on both retain and forget data. Let $\theta_b$ denote a checkpoint obtained before the model was exposed to the forget facts, and let $\theta$ denote the model being unlearned. For a prompt $x$, a partially masked response state $s$, and a masked position $i$, the model predicts
\begin{equation}
    p_{\theta,i}(v\mid x,s), \qquad v\in\V.
\end{equation}

The forget set is organized into facts $\{F_j\}_{j=1}^{J}$. Each fact has a set of semantic views, such as question--answer formulations, and a small fixed bank of MDLM states. The default state bank contains:
\begin{align}
    s^{\mathrm{recall}}(y)
    &= [\mathrm{MASK}]^{|y|}, \\
    s^{\mathrm{infill}}(y,K)
    &= \bigl(y_{\bar K},[\mathrm{MASK}]_{K}\bigr),
\end{align}
where $K$ is the fact-bearing span.

For each forget atom
\begin{equation}
    a=(x,s,i,y_i),
\end{equation}
we define a target-probability cap $\rho_a$. In counterfactual data-unlearning mode,
\begin{equation}
    \rho_a^{\mathrm{cf}}
    = p_{\theta_b,i}(y_i\mid x,s).
\end{equation}
If stricter knowledge erasure is required, use
\begin{equation}
    \rho_a
    = \min\!\left\{\rho_a^{\mathrm{cf}},\bar\rho_a\right\}.
\end{equation}

\section{One-sided Minimum-change Projection Loss}

Let
\begin{equation}
    p_a(\theta)
    := p_{\theta,i}(y_i\mid x,s).
\end{equation}
We require $p_a(\theta)\leq \rho_a$. The minimum-KL projection of a categorical distribution $p$ onto the set $\{q:q_{y_i}\leq\rho_a\}$ changes only the target probability and preserves the relative odds among all alternatives. Its scalar loss is
\begin{equation}
\label{eq:projection-loss}
    \ell_{\rho_a}\!\left(p_a(\theta)\right)
    =
    \begin{cases}
        0,
        & p_a(\theta)\leq\rho_a, \\
        \rho_a\log\dfrac{\rho_a}{p_a(\theta)}
        +(1-\rho_a)\log\dfrac{1-\rho_a}{1-p_a(\theta)},
        & p_a(\theta)>\rho_a.
    \end{cases}
\end{equation}
Equivalently, the active branch is
\begin{equation}
    \KL\!\left(
        \Ber(\rho_a)
        \;\middle\|\;
        \Ber(p_a(\theta))
    \right).
\end{equation}
If $m_a$ is the target-versus-rest log-odds and $p_a=\sigma(m_a)$, then on the active branch
\begin{equation}
    \frac{\partial \ell_{\rho_a}}{\partial m_a}
    = p_a-\rho_a.
\end{equation}
Thus, the gradient decreases smoothly to zero at the desired forgetting boundary and does not continue suppressing the token after the cap is met.

\section{Fact-wise Forgetting Constraints}

For fact $F_j$, let $\A_j$ be the set of all monitored atoms across supplied semantic views, recall states, infill states, and selected fact-bearing positions. Define
\begin{equation}
\label{eq:fact-constraint}
    \mathcal{F}_j(\theta)
    = \max_{a\in\A_j}
      \ell_{\rho_a}\!\left(p_a(\theta)\right).
\end{equation}
The constraint
\begin{equation}
    \mathcal{F}_j(\theta)\leq\delta_j
\end{equation}
requires every monitored atom of fact $j$ to satisfy the desired forgetting tolerance. Unlike an average forget loss, one easy fact cannot compensate for another fact that remains strongly memorized.

\section{Retain-behavior Objective}

The retain objective preserves the original MDLM distribution on retained examples:
\begin{equation}
\label{eq:retain-loss}
\begin{aligned}
    \mathcal{R}_{\mathrm{retain}}(\theta,\theta_0)
    = \mathbb{E}_{\substack{(x_r,y_r)\sim\Dr\\M\sim Q_r}}
    \left[
        \frac{1}{|M|}
        \sum_{i\in M}
        \KL\!\left(
            p_{\theta_0,i}(\cdot\mid x_r,s_M^{y_r})
            \;\middle\|\;
            p_{\theta,i}(\cdot\mid x_r,s_M^{y_r})
        \right)
    \right].
\end{aligned}
\end{equation}
Teacher logits from $\theta_0$ may be cached on a fixed retain-mask bank to remove the online teacher forward pass.

The constrained training problem is
\begin{equation}
\label{eq:main-objective}
\boxed{
\begin{aligned}
    \min_{\theta}\quad
    & \mathcal{R}_{\mathrm{retain}}(\theta,\theta_0) \\
    \text{subject to}\quad
    & \mathcal{F}_j(\theta)\leq\delta_j,
      \qquad j=1,\ldots,J.
\end{aligned}
}
\end{equation}

\section{Precomputation of Counterfactual Caps}

Define the target-versus-rest log-odds
\begin{equation}
    m_{\theta}(a)
    = z_{\theta,i}(y_i\mid x,s)
      - \log\sum_{v\neq y_i}\exp z_{\theta,i}(v\mid x,s).
\end{equation}
The acquired-evidence score is
\begin{equation}
    d(a)
    = \relu{m_{\theta_0}(a)-m_{\theta_b}(a)}.
\end{equation}
This score is a checkpoint-based attribution proxy. It identifies positions whose target evidence increased after exposure to the forget data.

For each fact, one may retain all non-prompt-copied target atoms, or select the shortest prefix of atoms sorted by $d(a)$ that covers a chosen fraction $\kappa$ of total acquired evidence.

\begin{algorithm}[t]
\caption{Counterfactual atom-bank and cap construction}
\label{alg:precompute}
\begin{algorithmic}[1]
\Require Original model $\theta_0$; pre-exposure checkpoint $\theta_b$; forget facts $\{F_j\}_{j=1}^{J}$; evidence coverage $\kappa\in(0,1]$; optional strict caps $\{\bar\rho_a\}$.
\Ensure Fact atom banks $\{\A_j\}$ and scalar caps $\{\rho_a\}$.
\For{$j=1,\ldots,J$}
    \State $\A_j^{\mathrm{cand}}\gets\varnothing$
    \ForAll{provided semantic views $(x,y)$ of fact $F_j$}
        \State Construct $s^{\mathrm{recall}}(y)$ and $s^{\mathrm{infill}}(y,K)$.
        \ForAll{state $s\in\{s^{\mathrm{recall}},s^{\mathrm{infill}}\}$}
            \State Run frozen forward passes of $\theta_0$ and $\theta_b$ on $(x,s)$.
            \ForAll{fact-bearing masked positions $i$ not copied from the prompt}
                \State $a\gets(x,s,i,y_i)$.
                \State $d(a)\gets\relu{m_{\theta_0}(a)-m_{\theta_b}(a)}$.
                \State $\rho_a^{\mathrm{cf}}\gets p_{\theta_b,i}(y_i\mid x,s)$.
                \If{strict-erasure mode is enabled}
                    \State $\rho_a\gets\min\{\rho_a^{\mathrm{cf}},\bar\rho_a\}$.
                \Else
                    \State $\rho_a\gets\rho_a^{\mathrm{cf}}$.
                \EndIf
                \State Add $(a,d(a),\rho_a)$ to $\A_j^{\mathrm{cand}}$.
            \EndFor
        \EndFor
    \EndFor
    \State Sort $\A_j^{\mathrm{cand}}$ by $d(a)$ in descending order.
    \State Let $\A_j$ be the shortest prefix satisfying
    \Statex \hspace{1.5em}$\displaystyle
        \sum_{a\in\A_j}d(a)
        \geq
        \kappa\sum_{a\in\A_j^{\mathrm{cand}}}d(a)$.
    \If{the total acquired evidence is zero}
        \State Set $\A_j\gets\A_j^{\mathrm{cand}}$ or mark $F_j$ as already counterfactual.
    \EndIf
\EndFor
\State \Return $\{\A_j\}_{j=1}^{J}$ and $\{\rho_a\}$.
\end{algorithmic}
\end{algorithm}

\section{Primal--Dual Training}

Introduce one nonnegative multiplier $\lambda_j$ for each fact. The Lagrangian is
\begin{equation}
\label{eq:lagrangian}
    \mathcal{L}(\theta,\bm\lambda)
    = \mathcal{R}_{\mathrm{retain}}(\theta,\theta_0)
      + \frac{1}{J}\sum_{j=1}^{J}
        \lambda_j\bigl(\mathcal{F}_j(\theta)-\delta_j\bigr),
    \qquad \lambda_j\geq0.
\end{equation}
The primal and dual updates are
\begin{align}
    \theta
    &\leftarrow
    \theta-
    \eta_{\theta}\nabla_{\theta}
    \mathcal{L}(\theta,\bm\lambda),
    \\
    \lambda_j
    &\leftarrow
    \relu{
        \lambda_j+
        \eta_{\lambda}\bigl(\mathcal{F}_j(\theta)-\delta_j\bigr)
    }.
\end{align}
Facts that already satisfy their caps receive no projection gradient. Facts that remain violated automatically receive stronger dual pressure.

\begin{algorithm}[t]
\caption{\FCP{} training}
\label{alg:training}
\begin{algorithmic}[1]
\Require Initial model $\theta_0$; retain data $\Dr$; atom banks $\{\A_j\}$; caps $\{\rho_a\}$; tolerances $\{\delta_j\}$; primal learning rate $\eta_{\theta}$; dual learning rate $\eta_{\lambda}$; feasibility-check period $T_{\mathrm{check}}$; patience $P$.
\Ensure Unlearned model $\theta$.
\State $\theta\gets\theta_0$.
\State Initialize $\lambda_j\geq0$ for all $j$.
\State $c_{\mathrm{feasible}}\gets0$.
\For{$t=1,2,\ldots,T_{\max}$}
    \State Sample a minibatch of fact indices $\mathcal{B}_f\subseteq\{1,\ldots,J\}$.
    \State Pack all required recall/infill states for $\{\A_j:j\in\mathcal{B}_f\}$ into one MDLM minibatch.
    \State Run one student forward pass and obtain $p_a(\theta)$ for every sampled atom.
    \ForAll{$j\in\mathcal{B}_f$}
        \State $\displaystyle
        \widehat{\mathcal{F}}_j
        \gets
        \max_{a\in\A_j}
        \ell_{\rho_a}\!\left(p_a(\theta)\right)$.
        \Comment{Eq.~\eqref{eq:fact-constraint}}
    \EndFor
    \State Sample a retain minibatch and a standard random mask pattern.
    \State Compute $\widehat{\mathcal{R}}_{\mathrm{retain}}$ using Eq.~\eqref{eq:retain-loss}.
    \State $\displaystyle
    \widehat{\mathcal{L}}
    \gets
    \widehat{\mathcal{R}}_{\mathrm{retain}}
    +\frac{1}{|\mathcal{B}_f|}
    \sum_{j\in\mathcal{B}_f}
    \lambda_j\bigl(\widehat{\mathcal{F}}_j-\delta_j\bigr)$.
    \State Update $\theta\gets\theta-\eta_{\theta}\nabla_{\theta}\widehat{\mathcal{L}}$.
    \ForAll{$j\in\mathcal{B}_f$}
        \State $\displaystyle
        \lambda_j
        \gets
        \relu{
            \lambda_j+
            \eta_{\lambda}
            \bigl(\stopgrad(\widehat{\mathcal{F}}_j)-\delta_j\bigr)
        }$.
    \EndFor
    \If{$t\bmod T_{\mathrm{check}}=0$}
        \State Evaluate every $\mathcal{F}_j(\theta)$ on the complete fixed atom bank.
        \If{$\mathcal{F}_j(\theta)\leq\delta_j$ for all $j$}
            \State $c_{\mathrm{feasible}}\gets c_{\mathrm{feasible}}+1$.
        \Else
            \State $c_{\mathrm{feasible}}\gets0$.
        \EndIf
        \If{$c_{\mathrm{feasible}}\geq P$}
            \State \textbf{break}
        \EndIf
    \EndIf
\EndFor
\State \Return $\theta$.
\end{algorithmic}
\end{algorithm}

\section{Optional Strict-cap Calibration}

Suppose a fact-bearing span contains $q$ selected positions and the desired exact-span probability at a fixed MDLM state is at most $\varepsilon_j$. A uniform per-position cap can be chosen as
\begin{equation}
    \bar\rho_a=\varepsilon_j^{1/q}.
\end{equation}
Under the standard factorized MDLM prediction head,
\begin{equation}
    p_{\theta}(y_K\mid x,s)
    = \prod_{i\in K}p_{\theta,i}(y_i\mid x,s).
\end{equation}
Therefore, if every selected position satisfies $p_{\theta,i}(y_i\mid x,s)\leq\varepsilon_j^{1/q}$, then
\begin{equation}
    p_{\theta}(y_K\mid x,s)\leq\varepsilon_j.
\end{equation}
This is a fixed-state exact-span guarantee; it is not a certificate for every iterative decoding trajectory or unseen semantic formulation.

\section{Implementation Notes}

\paragraph{No adversarial inner loop.}
The training algorithm uses a fixed recall/infill state bank. It requires no mask restarts, embedding-gradient attack, discrete swaps, or repeated candidate evaluation.

\paragraph{Fact batching.}
Facts should be sampled uniformly rather than sampling individual atoms uniformly. This prevents facts with many views or long answers from dominating training.

\paragraph{Hard maximum.}
The core method uses a hard maximum over atoms of each fact. A top-$s$ mean or log-sum-exp relaxation may improve optimization stability, but should be treated as an ablation because it weakens the exact per-fact interpretation.

\paragraph{Teacher caching.}
The pre-exposure model $\theta_b$ is used only during offline cap construction. The original retain teacher $\theta_0$ can also be removed from the online loop by caching its logits on a fixed retain-mask bank.

\paragraph{Stopping criterion.}
Training should stop according to complete-bank feasibility rather than a fixed epoch count. Report the fact pass rate
\begin{equation}
    \operatorname{FactPassRate}(\theta)
    = \frac{1}{J}
      \sum_{j=1}^{J}
      \mathbf{1}\!\left[\mathcal{F}_j(\theta)\leq\delta_j\right],
\end{equation}
together with the maximum and upper quantiles of $\{\mathcal{F}_j(\theta)\}_{j=1}^{J}$.

\paragraph{Scope of the claim.}
The algorithm enforces forgetting over the declared semantic views, token positions, and MDLM state bank. Held-out paraphrases, alternative mask geometries, decoding schedules, membership-inference probes, and relearning attacks should be evaluated separately.

\end{document}
